% !TEX program = xelatex
\documentclass{misisarticle}
\grnti{49.33.29; 28.17.19}
\asjc{1705; 2208; 2604}
\oecd{1.02; 2.02}

\title{Экспериментальная оценка децентрализованной передачи сообщений в LoRa-mesh сетях для оффлайн-коммуникаций}

\misisauthor{Соловьев Вадим Дмитриевич}
  {кафедра инженерной кибернетики\par НИТУ «МИСИС», БПМ-23-ПО-1}
  {Москва, Россия}
  {}


\setabstract{Рассматривается передача коротких сообщений в децентрализованной LoRa mesh-сети без инфраструктуры оператора. Цель исследования --- разделить вычислительные затраты поиска пути и затраты протокола, который обнаруживает и поддерживает маршруты в ограниченном радиоканале. Описана двухуровневая методика: сравнение шести алгоритмов поиска пути на синтетических графах и дискретно-событийное моделирование пяти стратегий передачи сообщений. Уровни связывает стоимость канала, выраженная через вероятность приёма и время передачи кадра. Модель учитывает потери, коллизии, полудуплексный режим, ограничение рабочего цикла и расход энергии. Заданы метрики доставки, задержки, служебного эфира и энергоэффективности, а также порядок сравнения на одинаковых реализациях сети. Обсуждаются ограничения модели и условия последующей проверки на данных действующей сети Meshtastic.}

\setkeywords{LoRa mesh-сеть, управляемое лавинное распространение, AODV/OLSR, рабочий цикл, ETX, дискретно-событийное моделирование}

\hypersetup{
  pdftitle={Экспериментальная оценка децентрализованной передачи сообщений в LoRa-mesh сетях для оффлайн-коммуникаций},
  pdfauthor={Соловьев Вадим Дмитриевич},
  pdfsubject={Научное письмо. ПЗ-3. Постановка задачи и методы исследования}
}

\begin{document}
\maketitle
\section{Введение}

Оффлайн-коммуникации обеспечивают обмен сообщениями между пользователями при отсутствии доступа к сети оператора или интернету. Такая связь востребована на удалённых территориях, в экспедициях и при нарушении работы инфраструктуры. Для коротких текстовых сообщений существенны надёжность доставки и автономность узлов, тогда как высокая скорость передачи не всегда необходима. Практическую возможность маломощной системы обмена сообщениями на основе LoRa рассматривают Cardenas и соавторы~\cite{cardenas}.

LoRa задаёт физический способ передачи, но сама по себе не определяет маршрутизацию в многошаговой сети. В mesh-сети промежуточные устройства ретранслируют кадры, увеличивая достижимую область связи. Каждый дополнительный переход при этом расходует эфир и энергию. Увеличение коэффициента расширения сигнала SF изменяет чувствительность и длительность передачи, поэтому выигрыш в дальности необходимо сопоставлять с ростом занятости канала~\cite{semtech}.

Алгоритм поиска пути и протокол маршрутизации решают разные задачи. Алгоритм получает готовый граф и вычисляет путь. Протокол должен сначала обнаружить соседей, получить сведения о каналах и поддерживать их при отказах узлов. Служебные кадры используют ту же среду, что и сообщения пользователей. Поэтому малое время вычисления кратчайшего пути ещё не гарантирует высокую вероятность доставки.

Цель работы --- разработать методику оценки децентрализованной передачи сообщений, которая позволяет отдельно исследовать вычислительный уровень и сетевые затраты. Для этого используются два связанных стенда: алгоритмический и дискретно-событийный. Первый сравнивает поиск пути при известных весах, второй --- стратегии передачи при потерях и конкуренции за эфир. Связь между стендами задаётся физически интерпретируемой стоимостью ребра.

Вклад работы состоит в организации такого сравнения на согласованных входных данных и в явном разграничении оптимального одноадресного маршрута и стратегии с несколькими потенциальными ретрансляторами. Настоящая версия статьи содержит постановку задачи и описание методов. Численные результаты и итоговые рекомендации по выбору протокола в неё не включены.

\section{Анализ работ}

\subsection{Передача сообщений и многошаговые LoRa-сети}

Cardenas и соавторы~\cite{cardenas} описывают недорогую маломощную систему обмена сообщениями на LoRa. Эта работа близка по прикладной задаче: связь требуется пользователям, которым недоступна обычная инфраструктура. Для рассматриваемого исследования она задаёт сценарий коротких сообщений, однако не заменяет отдельной оценки стоимости маршрутизации.

Durand и Booysen~\cite{durand} исследуют mesh-топологию LoRa с помощью модели ns-3. Они анализируют доставку к шлюзу и энергопотребление при различных сценариях. Работа показывает необходимость совместного рассмотрения покрытия, ретрансляций и энергии. В отличие от схемы доставки к центральному шлюзу, здесь рассматривается обмен между парами равноправных узлов.

Udugampola и соавторы~\cite{udugampola} предлагают маршрутизацию с учётом положения и энергии узлов для подземных LoRa-сетей. В ней используются начальная фаза определения относительного положения, резервные ретрансляторы и переключение маршрутов. Результаты такой схемы зависят от геометрии и нагрузки подземной сети. Они мотивируют сравнение механизмов при различных условиях, но не дают универсального ранжирования протоколов для наземной оффлайн-связи.

Bor и соавторы~\cite{bor} экспериментально исследуют поведение LoRa и используют полученные зависимости в LoRaSim. Их подход важен для выбора модели потерь и обработки пересекающихся передач. Параметры распространения, измеренные в конкретной среде, при переносе в другую местность требуют калибровки.

\subsection{Маршрутизация и стоимость канала}

AODV обнаруживает маршрут по запросу и поддерживает его по мере необходимости~\cite{aodv}. OLSR заранее распространяет сведения о соседях и топологии, применяя многоточечные ретрансляторы MPR для сокращения служебного распространения~\cite{olsr}. Эти подходы позволяют сравнить затраты реактивного обнаружения и проактивного поддержания маршрутов.

Альтернативой служит управляемое лавинное распространение. Документация Meshtastic описывает задержку ретрансляции с учётом качества принятого сигнала и подавление части повторных передач после приёма дубликата~\cite{meshtastic}. В данной работе реализована упрощённая схема этого типа. Она воспроизводит выбранные механизмы распространения, но не весь стек и не все режимы прошивки Meshtastic.

Метрика ETX, предложенная De Couto и соавторами~\cite{etx}, учитывает ожидаемое число передач с повторениями и позволяет оценивать каналы иначе, чем по числу переходов. В классической постановке существенны вероятности доставки данных и подтверждения. В разработанном стенде используется упрощённая стоимость, пропорциональная времени передачи и обратно пропорциональная вероятности приёма кадра. Её ограничения при наличии подтверждений рассматриваются ниже.

Классические алгоритмы поиска кратчайших путей описаны, в частности, в~\cite{cormen}. Их сравнение характеризует вычисления на готовом графе. Для оценки LoRa mesh необходимо дополнительно учитывать, каким образом этот граф становится известен узлам и какой объём эфира расходуется на его поддержание.

Таким образом, в работе сопоставляются два вида затрат: обработка графа в памяти и получение сетевого состояния посредством радиопередач. Методика не предполагает, что самый быстрый алгоритм автоматически определяет лучший протокол, и проверяет эти уровни отдельно.

\section{Постановка задачи}

\subsection{Входные данные и искомое решение}

Сеть представляется графом $G(t)=(V(t),E(t))$. Вершины соответствуют узлам, способным передавать и ретранслировать сообщения. Набор доступных каналов зависит от размещения устройств, параметров радиосвязи и состояния узлов. В имитационной модели координаты фиксированы в пределах прогона, а динамика задаётся отказами и восстановлением.

Входными данными являются координаты узлов, параметры LoRa, модель потерь, расписание прикладных сообщений и событий доступности. Каждое сообщение задаётся источником $s$, получателем $d$, временем создания и длиной полезной нагрузки. Общий канал полудуплексный: узел не может одновременно передавать и принимать.

Для алгоритмического уровня требуется найти путь $P$ между $s$ и $d$ с минимальной суммой неотрицательных весов:
\begin{equation}
  P^*=\arg\min_{P:s\leadsto d}\sum_{(u,v)\in P}w_{uv}.
  \label{eq:path}
\end{equation}
Наряду со стоимостью результата измеряются время вычисления, число посещённых вершин и число обработанных рёбер. Если пути нет, это фиксируется как отдельный исход. BFS минимизирует число переходов и используется как базовый метод, который не учитывает величину весов.

Для сетевого уровня требуется сопоставить стратегии передачи по доле доставленных сообщений, задержке, затратам эфира и энергии. Основной вопрос исследования: в каких условиях различия в доставке связаны с обслуживанием маршрутов и использованием радиоканала, даже если задача поиска пути на готовом графе решается эффективно?

\subsection{Метрики}

Коэффициент доставки определяется по уникальным прикладным сообщениям:
\begin{equation}
  \mathrm{PDR}=\frac{N_{\mathrm{del}}}{N_{\mathrm{gen}}},
\end{equation}
где $N_{\mathrm{gen}}$ --- число созданных сообщений, а $N_{\mathrm{del}}$ --- число сообщений, принятых конечным получателем. Повторные копии не увеличивают числитель. Задержка считается только для доставленных сообщений от создания до первого приёма адресатом. Поэтому задержку необходимо анализировать совместно с PDR: низкая задержка у небольшой доли успешных сообщений не означает высокое качество связи.

Служебные расходы характеризуются долей суммарного времени передач:
\begin{equation}
  H=\frac{T_{\mathrm{ctrl}}}{T_{\mathrm{ctrl}}+T_{\mathrm{data}}}.
\end{equation}
Это отношение сумм длительностей кадров, а не доля неперекрывающегося физического времени наблюдения. Кроме $H$ учитывается число передач на одно доставленное сообщение. У лавинной схемы может отсутствовать отдельный служебный трафик маршрутизации, но повторная передача данных всё равно расходует эфир.

Энергоэффективность задаётся как
\begin{equation}
  E_{\mathrm{bit}}=\frac{E_{\mathrm{net}}}{8L N_{\mathrm{del}}},
\end{equation}
где $E_{\mathrm{net}}$ --- суммарная энергия узлов, а $L$ --- длина полезной нагрузки в байтах. При отсутствии доставленных сообщений показатель не определён. Дополнительно рассматриваются полезная скорость доставки, время до истощения первой батареи и изменение PDR при отказах.

Окончательный выбор стратегии предполагает компромисс между надёжностью, задержкой и энергией. Единственная взвешенная сумма этих показателей не вводится: предпочтение зависит от требований конкретного сценария оффлайн-связи.

\section{Описание методов}

\subsection{Алгоритмическая модель}

Первый стенд сравнивает BFS, алгоритм Дейкстры, двунаправленный алгоритм Дейкстры, A*, алгоритмы Беллмана--Форда и Флойда--Уоршелла. Набор включает 12 типов топологий: линию, кольцо, звезду, решётку, случайный граф, плотную mesh-сеть, две области с узким соединением, безмасштабную сеть, малый мир, дерево, fat-tree и dragonfly. Последние топологии используются для исследования путевой избыточности и не интерпретируются как геометрические модели размещения LoRa-устройств.

Для каждой топологии подготовлено 30 графов, всего 360. Размер зависит от семейства: для линии, кольца и звезды генератор задаёт от 20 до 300 вершин, для решётки --- от 4 до 20 строк и столбцов, для двоичного дерева --- высоту от 3 до 7. Исходные веса --- целые числа от 1 до 50. Все алгоритмы получают один и тот же граф и пару вершин. Для взвешенных методов проверяется совпадение стоимости пути с опорным решением. У A* корректность зависит от допустимости эвристики; отсутствие информативной геометрической оценки ограничивает возможное ускорение.

Время выполнения характеризует конкретную программную реализацию и вычислительную машину. Число посещённых вершин и обработанных рёбер даёт дополнительную характеристику объёма работы. Измерения на рабочей станции не являются прямым измерением задержки или расхода энергии микроконтроллера. Алгоритм Флойда--Уоршелла решает более широкую задачу всех пар, что учитывается при интерпретации его затрат на один запрос.

\subsection{Дискретно-событийная модель радиоканала}

Второй стенд моделирует создание сообщений, ожидание доступа к среде, передачу и приём кадров, ретрансляцию и изменение состояния узлов. События обрабатываются в порядке модельного времени. Общая реализация физического уровня и доступа к среде используется для всех сравниваемых стратегий.

Длительность кадра $T_{\mathrm{air}}$ вычисляется по формуле Semtech~\cite{semtech}. Длительность символа равна
\begin{equation}
  T_{\mathrm{sym}}=\frac{2^{SF}}{BW}.
\end{equation}
Полное время включает преамбулу, заголовок, полезную нагрузку, проверочную сумму и избыточность кодирования. Поэтому при одинаковом числе байтов увеличение SF увеличивает занятость канала.

Потери распространения задаются моделью
\begin{equation}
\begin{split}
  PL(d)&=PL(d_0)+10n\log_{10}(d/d_0)+X,\\
  X&\sim\mathcal{N}(0,\sigma^2).
\end{split}
\end{equation}
В базовой конфигурации используются $PL(d_0)=127{,}41$~дБ при $d_0=40$~м, $n=2{,}08$ и $\sigma=3{,}57$~дБ, принятые по~\cite{bor}. Это параметры выбранной среды, а не универсальная характеристика городского канала.

По мощности передатчика, усилению антенн и потерям определяется RSSI, затем отношение сигнал/шум SNR. Чувствительность и порог SNR зависят от SF. Вблизи порога декодирование описывается вероятностной зависимостью. Для совпадающих во времени передач проверяется возможность захвата более сильного сигнала с порогом 6~дБ. Такой порог является упрощением модели интерференции.

Доступ к среде включает обнаружение занятости, случайную отсрочку и ограничение числа попыток. Рабочий цикл в базовом сценарии равен 1~\% и учитывается в скользящем окне 3600~с. Кадр, превышающий бюджет узла, отбрасывается. Значение 1~\% задаёт исследуемый сценарий; требования конкретного частотного поддиапазона и региона должны устанавливаться отдельно.

Для одноадресных кадров моделируются подтверждения и до трёх повторных передач. Широковещательные кадры не подтверждаются. Энергия рассчитывается по состояниям передачи, приёма, ожидания и сна:
\begin{equation}
  E_i=U\sum_k I_k t_{ik}.
\end{equation}
Используются постоянные токи состояний. Модель не воспроизводит нелинейный разряд батареи и индивидуальные характеристики всех аппаратных платформ.

\begin{table}[t]
\caption{Базовые параметры сетевого стенда}
\label{tab:params}
\centering
\small
\begin{tabular}{@{}p{0.58\columnwidth}r@{}}
\toprule
Параметр & Значение \\
\midrule
Число узлов & 25 \\
Коэффициент расширения & SF9 \\
Полоса & 125 кГц \\
Скорость кодирования & 4/5 \\
Мощность передатчика & 14 дБм \\
Полезная нагрузка / заголовок & 32 / 16 байт \\
Интенсивность для всей сети & 6 сообщений/мин \\
Основное модельное время & 1200 с \\
Завершение оставшихся передач & 120 с \\
Рабочий цикл & 1\,\% \\
Ёмкость батареи & 2000 мА$\cdot$ч \\
\bottomrule
\end{tabular}
\end{table}

\subsection{Сравниваемые стратегии}

\textit{AODV.} Используется модель обнаружения пути сообщениями запроса и ответа. Она позволяет исследовать стоимость поиска и восстановления маршрута. Это упрощённая реализация механизмов AODV, а не проверка полного соответствия RFC~3561.

\textit{OLSR.} Моделируются периодические сообщения о соседях и топологии и распространение с использованием MPR. Цель --- оценить стоимость поддержания маршрутов при дефиците эфира. Полнота промышленной реализации RFC~3626 не предполагается.

\textit{ManagedFlood.} Узел ретранслирует ранее не встречавшееся сообщение с ограничением числа переходов. Время ожидания зависит от SNR: меньший запас сигнала даёт более раннюю попытку. Если до неё принят дубликат, запланированная передача отменяется. В базовой реализации предел равен восьми переходам.

\textit{EpidemicGossip.} Моделируется вероятностное лавинное распространение: после первого приёма сообщение передаётся с вероятностью 0,8, предел --- десять переходов. В отличие от ManagedFlood подавление ожидающей передачи по дубликату отсутствует. Здесь термин обозначает gossip-схему и не подразумевает полноценный DTN-протокол с обменом буферами при встречах узлов.

\textit{Ideal.} Эталон мгновенно получает доступный граф и выбирает путь минимальной стоимости без служебного обнаружения. Передачи по выбранному маршруту подчиняются тем же ограничениям канала. Эталон относится к одноадресной передаче по одному пути и не является верхней границей PDR для произвольного способа распространения.

\subsection{Связь уровней через физические веса}

Алгоритмический стенд повторно применяется к графам, построенным сетевой моделью. Для ребра задаётся стоимость
\begin{equation}
  w_{uv}=\frac{T_{\mathrm{air}}}{p_{uv}},
  \label{eq:weight}
\end{equation}
где $p_{uv}$ --- вероятность приёма кадра при среднем состоянии линии без текущих коллизий. Стоимость имеет размерность времени и характеризует ожидаемые затраты на успешную передачу в упрощённой модели независимых попыток. При ограниченном числе повторений, наличии подтверждений и конкуренции за эфир она служит маршрутной оценкой, а не точным временем доставки.

В реализации различаются граф физически слышимых соседей и граф каналов для маршрутизации. Во второй включаются линии с $p_{uv}\geq0{,}7$; по нему строит путь Ideal. Порог и множество допустимых каналов влияют на сравнение: лавинное распространение может использовать приём вне заранее выбранного одноадресного пути.

Для проверки переноса весов используются три типа размещения, три значения seed и до 12 пар вершин на граф. На этих графах повторяются алгоритмы первого стенда и проверяется стоимость решения~\eqref{eq:path}. Эта проверка связывает уровни задачи, но сама по себе не доказывает превосходство какой-либо стратегии передачи.

\subsection{Сценарии и статистическая обработка}

Основное сравнение включает статическую сеть, случайные отказы с восстановлением и одновременное отключение 35~\% узлов. В потоке сообщений интервалы поступления имеют экспоненциальное распределение. Пары отправителя и получателя выбираются среди связанных в исходном маршрутном графе узлов. Поэтому PDR характеризует передачу для этого набора пар, а не покрытие всех возможных пар на территории.

Для основного сравнения задаются десять seed от 42 до 51. При одинаковом seed все стратегии получают одинаковое размещение, поток прикладных сообщений и расписание внешних отказов. Радиособытия и потребление случайных чисел после начала работы могут различаться из-за поведения протоколов. Для каждого показателя рассчитываются среднее и приближённый 95-процентный доверительный интервал $\bar{x}\pm1{,}96s/\sqrt{r}$, где $s$ --- выборочное стандартное отклонение, а $r$ --- число независимых прогонов. При малом числе повторений точность нормального приближения ограничена. Наблюдения внутри одной сети не считаются независимыми повторениями эксперимента.

Дополнительные серии изменяют число узлов от 15 до 80, SF от 7 до 12, интенсивность нагрузки и размещение: равномерное, кластерное, коридорное и регулярное. Число повторений в этих сериях определяется отдельно в настройках запуска. При изменении SF генератор размещения использует соответствующую расчётную дальность, поэтому такую серию следует интерпретировать как сравнение нормированных развёртываний, а не как изменение SF на одной фиксированной карте.

Результаты должны представляться вместе с конфигурацией, горизонтом моделирования и числом повторений. Сравнение PDR сопровождается задержкой, служебным эфиром и числом передач на доставку. Это позволяет отличать уменьшение накладных расходов от перераспределения затрат в повторные передачи данных.

\subsection{Ограничения и последующая проверка}

Ранжирование относится к заданным моделям, интервалам служебных сообщений и сценариям. Оно не переносится автоматически на все реализации AODV, OLSR или прошивки Meshtastic. Одноканальная модель с одинаковыми параметрами устройств не описывает всю неоднородность городской сети. Измерения на компьютере также не устанавливают вычислительную стоимость на встроенной платформе.

Следующий этап исследования посвящён физическому уровню действующей сети Meshtastic. Предполагается сопоставить модель с наблюдаемыми характеристиками каналов и затем проверить прогноз собственными узлами известной конфигурации. Данные шлюзов содержат только наблюдаемые ими пакеты, поэтому частоту приёма через шлюз нельзя напрямую приравнять к вероятности прямой доставки между всеми узлами. Параметры московской сети и её радиопресет должны задаваться отдельно от базового сценария таблицы~\ref{tab:params}. Завершённая натурная проверка в данной статье не заявляется.

\begin{thebibliography}{99}

\bibitem{cardenas}
Cardenas A.~M., Nakamura Pinto M.~K., Pietrosemoli E., Zennaro M., Rainone M., Manzoni P. A Low-Cost and Low-Power Messaging System Based on the LoRa Wireless Technology // Mobile Networks and Applications. 2020. Vol.~25, no.~3. P.~961--968. DOI: \href{https://doi.org/10.1007/s11036-019-01235-5}{10.1007/s11036-019-01235-5}.

\bibitem{semtech}
Semtech Corporation. SX1276/77/78/79: 137 MHz to 1020 MHz Low Power Long Range Transceiver. Datasheet. Rev.~7. 2020. URL: \url{https://www.semtech.com/products/wireless-rf/lora-connect/sx1276} (дата обращения: 08.11.2026).

\bibitem{durand}
Durand T.~G., Booysen M.~J. Performance Evaluation of a Mesh-Topology LoRa Network // Sensors. 2025. Vol.~25, no.~5. Art.~1602. DOI: \href{https://doi.org/10.3390/s25051602}{10.3390/s25051602}.

\bibitem{udugampola}
Udugampola N., Ai X., Li B., Gong H., Seneviratne A. A Position- and Energy-Aware Routing Strategy for Subterranean LoRa Mesh Networks. 2025. arXiv:2510.03714. DOI: \href{https://doi.org/10.48550/arXiv.2510.03714}{10.48550/arXiv.2510.03714}.

\bibitem{bor}
Bor M., Roedig U., Voigt T., Alonso J. Do LoRa Low-Power Wide-Area Networks Scale? // Proceedings of the 19th ACM International Conference on Modeling, Analysis and Simulation of Wireless and Mobile Systems. 2016. P.~59--67. DOI: \href{https://doi.org/10.1145/2988287.2989163}{10.1145/2988287.2989163}.

\bibitem{aodv}
Perkins C., Belding-Royer E., Das S. Ad hoc On-Demand Distance Vector (AODV) Routing. RFC~3561. IETF, 2003. DOI: \href{https://doi.org/10.17487/RFC3561}{10.17487/RFC3561}.

\bibitem{olsr}
Clausen T., Jacquet P. Optimized Link State Routing Protocol (OLSR). RFC~3626. IETF, 2003. DOI: \href{https://doi.org/10.17487/RFC3626}{10.17487/RFC3626}.

\bibitem{meshtastic}
Meshtastic Project. Mesh Broadcast Algorithm. Документация, версия 2.8. URL: \url{https://meshtastic.org/docs/overview/mesh-algo/} (дата обращения: 08.11.2026).

\bibitem{etx}
De Couto D.~S.~J., Aguayo D., Bicket J., Morris R. A High-Throughput Path Metric for Multi-Hop Wireless Routing // Proceedings of the 9th Annual International Conference on Mobile Computing and Networking. 2003. P.~134--146. DOI: \href{https://doi.org/10.1145/938985.939000}{10.1145/938985.939000}.

\bibitem{cormen}
Cormen T.~H., Leiserson C.~E., Rivest R.~L., Stein C. Introduction to Algorithms. 3rd~ed. Cambridge, MA: MIT Press, 2009. 1312~p.

\end{thebibliography}

\end{document}
